% Source: https://artofproblemsolving.com/community/c3h2671362_median_is_so_confusing_now_
% Stable identity: aops-topic-2671362-q01
\textbf{(a)} For $1,2,3$ we have $n=3$ (odd), and the formula gives $\dfrac{X[1]+X[2]}{2}=\dfrac{1+2}{2}=1.5$, but the median (the middle value) is $2$. For $1,2,3,4$ we have $n=4$ (even), and the formula gives $X[2]=2$, but the median (the average of the two middle values) is $\dfrac{2+3}{2}=2.5$.

\textbf{(b)} The median is the middle value when $n$ is odd and the average of the two middle values when $n$ is even:
\[\mathrm{Med}(X)=\begin{cases}X\left[\dfrac{n+1}{2}\right]&\text{if }n\text{ is odd,}\\[2ex]\dfrac{X\left[\dfrac n2\right]+X\left[\dfrac n2+1\right]}{2}&\text{if }n\text{ is even.}\end{cases}\]
For odd $n$ there are $\dfrac{n-1}{2}$ values on each side of $X\left[\dfrac{n+1}{2}\right]$; for even $n$ there are $\dfrac n2-1$ values on each side of the middle pair. So the website's formula has the two cases mixed up (and its indices off by one). The median is indeed just ``the middle number''.
